\chapter*{Vorwort}
\addcontentsline{toc}{chapter}{Vorwort}

\section*{Worum es geht}

\begin{merksatz}
\boxtitel{Ein Molekül, vier Symmetrien, drei Töne}
Das Wassermolekül hat vier Symmetrien: nichts tun, um $180^{\circ}$ drehen, an der Molekülebene spiegeln, an der Ebene senkrecht dazu spiegeln. Es hat drei Atome, also neun Freiheitsgrade der Bewegung; drei davon sind Verschiebungen des ganzen Moleküls, drei Drehungen, und drei sind Schwingungen --- die Töne, die man im Infrarotspektrum sieht.

\begin{center}
{\small
\renewcommand{\arraystretch}{1.5}
\begin{tabular}{@{}P{3.4cm}P{5.2cm}@{}}
\toprule
\textbf{Was man weiß} & die Symmetriegruppe (vier Elemente) und die Zahl der Atome\\
\textbf{Was man ausrechnet} & ohne eine einzige Kraft zu kennen: Die drei Schwingungen zerfallen in zwei symmetrische ($A_1$) und eine antisymmetrische ($B_2$); alle drei sind im Infrarot sichtbar; keine ist entartet.\\
\textbf{Was man misst} & $3657$, $1595$ und $3756\ \mathrm{cm}^{-1}$ --- zwei Banden, die sich bei der Drehung nicht ändern, eine, die das Vorzeichen wechselt.\\
\bottomrule
\end{tabular}}
\end{center}

\begin{center}
\begin{minipage}{0.9\textwidth}
\textbf{Die Symmetrie allein sagt, wie viele Schwingungen welcher Art ein Molekül hat, welche man sehen kann und welche zusammenfallen. Die Frequenzen sagt sie nicht --- aber alles, was an einem Spektrum Struktur ist, folgt aus einer Gruppe mit vier Elementen und einer Tafel mit sechzehn Einträgen.}
\end{minipage}
\end{center}

\begin{center}
\begin{minipage}{0.88\textwidth}
\kursiv{\textbf{Darstellungstheorie ist die Lehre davon, wie eine Gruppe auf einem Vektorraum wirkt} --- als Matrizen, die man multiplizieren kann --- und was sich daraus über den Raum ergibt, ohne ihn im Einzelnen zu kennen. Der Vektorraum ist hier der Raum der Bewegungen eines Moleküls, später der Raum seiner Elektronenzustände, der Funktionen auf einer Kugel, der Zustände eines Elektrons mit Spin, der Quarks in einem Proton. Immer ist die Frage dieselbe: In welche unzerlegbaren Stücke zerfällt der Raum unter der Symmetrie? Und immer ist die Antwort eine Tafel von Zahlen, die Charaktertafel, die man für jede endliche Gruppe berechnen kann. Der Leitfall dieses Bandes sind Moleküle --- Wasser, Ammoniak, Benzol ---, weil an ihnen die Theorie etwas vorhersagt, das man im Labor messen kann.}
\end{minipage}
\end{center}
\end{merksatz}

\section*{Drei Überraschungen}

\begin{merksatz}
\boxtitel{Was dieser Band zeigt}
\begin{center}
{\small
\renewcommand{\arraystretch}{1.5}
\begin{tabular}{@{}P{0.5cm}P{8.2cm}@{}}
\toprule
$1$ & Jede endliche Gruppe hat genau so viele unzerlegbare Darstellungen wie Konjugationsklassen, und die Summe der Quadrate ihrer Dimensionen ist die Gruppenordnung. Für die Symmetrien eines Dreiecks: $1^{2} + 1^{2} + 2^{2} = 6$. \kursiv{(Kapitel 4 und 5)}\\
$2$ & Die Charaktere --- die Spuren der Matrizen --- bestimmen die Darstellung vollständig, und sie bilden ein Orthonormalsystem. Eine Darstellung zu zerlegen ist ein Skalarprodukt. \kursiv{(Kapitel 4)}\\
$3$ & Die Kugelfunktionen, die Orbitale $s, p, d, f$ der Chemie, sind nichts als die unzerlegbaren Darstellungen der Drehgruppe --- und ihre Dimensionen $1, 3, 5, 7$ sind der Grund für den Aufbau des Periodensystems. \kursiv{(Kapitel 10)}\\
\bottomrule
\end{tabular}}
\end{center}

\begin{center}
\begin{minipage}{0.88\textwidth}
\kursiv{Die erste und zweite Überraschung sind Sätze von Frobenius aus den Jahren 1896 bis 1899, die dritte ist Wigners und Weyls Einsicht der zwanziger Jahre, die aus der Darstellungstheorie ein Werkzeug der Quantenmechanik machte. Alle drei werden in diesem Band bewiesen --- die ersten beiden vollständig, die dritte bis auf die Analysis der Drehgruppe, die als Auslassung benannt ist.}
\end{minipage}
\end{center}
\end{merksatz}

\section*{Der Aufbau}

\begin{merksatz}
\boxtitel{Die sechs Teile}
\begin{center}
{\small
\renewcommand{\arraystretch}{1.5}
\begin{tabular}{@{}P{1.4cm}P{4.2cm}P{2.4cm}@{}}
\toprule
\textbf{Teil} & \textbf{Thema} & \textbf{Kapitel}\\
\midrule
I & Gruppen und ihre Darstellungen & $1$--$3$\\
II & Charaktere & $4$--$6$\\
III & Konstruktionen & $7$--$9$\\
IV & Kontinuierliche Gruppen & $10$--$12$\\
V & Moleküle, Kristalle, Teilchen & $13$--$14$\\
VI & Stand & $15$\\
\bottomrule
\end{tabular}}
\end{center}

\schritt{Was vorausgesetzt wird}
Lineare Algebra mit Eigenwerten, Skalarprodukten und Matrizen (Band 2); Gruppen, Untergruppen und Nebenklassen in dem Umfang von Band 11 (Algebra); die komplexen Zahlen (Band 12). Für Teil IV ist etwas Analysis (Band 3, Band 13) hilfreich, für Kapitel 14 die Quantenmechanik so weit, wie sie in Band 20 der Informatik-Reihe oder in der Physik-Reihe erscheint --- aber nicht nötig: Der Band sagt an jeder Stelle, was er braucht. Alles andere wird hier entwickelt, und wo ein Beweis übergangen wird, steht es dabei.

\begin{center}
\begin{minipage}{0.88\textwidth}
\kursiv{\textbf{Jedes Kapitel endet mit einem Bruchfall}: einem Beispiel, an dem die Methode des Kapitels versagt --- die Gruppe, die fast eine Symmetriegruppe ist; der Charakter, der zwei Darstellungen nicht unterscheidet, weil die Charakteristik des Körpers nicht null ist; die Auswahlregel, die ein Übergang bricht. Und jedes Kapitel hat drei Übungsaufgaben mit vollständiger Lösung und Probe. Alle Charaktertafeln und Zerlegungen des Bandes wurden mit einem Programm von achtzig Zeilen gerechnet, das Gruppen aus Matrizen erzeugt, ihre Klassen bestimmt und Charaktere zerlegt; es liegt den Quellen bei.}
\end{minipage}
\end{center}
\end{merksatz}
